DGFiP, REI 2024 et DVF 2023-2024 ; Insee, recensement 2022, comptes de patrimoine et Insee Analyses n° 91 ; calculs de l'auteur.\par Note : taxe foncière des logements des particuliers rapportée à leur valeur de marché, par décile de logements classés selon le prix au m² de leur commune (France métropolitaine hors Alsace-Moselle). Les trois scénarios remplacent la moitié des droits de mutation (5,7 Md\euro{}) par de la taxe foncière, à recette constante. « Deux taux » : taux sur le terrain égal à quatre fois le taux sur le bâti.
